% Options for packages loaded elsewhere
% Options for packages loaded elsewhere
\PassOptionsToPackage{unicode}{hyperref}
\PassOptionsToPackage{hyphens}{url}
\PassOptionsToPackage{dvipsnames,svgnames,x11names}{xcolor}
%
\documentclass[
  korean,
  11pt,
  a4paper,
  DIV=11,
  numbers=noendperiod]{scrreprt}
\usepackage{xcolor}
\usepackage{amsmath,amssymb}
\setcounter{secnumdepth}{5}
\usepackage{iftex}
\ifPDFTeX
  \usepackage[T1]{fontenc}
  \usepackage[utf8]{inputenc}
  \usepackage{textcomp} % provide euro and other symbols
\else % if luatex or xetex
  \usepackage{unicode-math} % this also loads fontspec
  \defaultfontfeatures{Scale=MatchLowercase}
  \defaultfontfeatures[\rmfamily]{Ligatures=TeX,Scale=1}
\fi
\usepackage{lmodern}
\ifPDFTeX\else
  % xetex/luatex font selection
\fi
% Use upquote if available, for straight quotes in verbatim environments
\IfFileExists{upquote.sty}{\usepackage{upquote}}{}
\IfFileExists{microtype.sty}{% use microtype if available
  \usepackage[]{microtype}
  \UseMicrotypeSet[protrusion]{basicmath} % disable protrusion for tt fonts
}{}
\makeatletter
\@ifundefined{KOMAClassName}{% if non-KOMA class
  \IfFileExists{parskip.sty}{%
    \usepackage{parskip}
  }{% else
    \setlength{\parindent}{0pt}
    \setlength{\parskip}{6pt plus 2pt minus 1pt}}
}{% if KOMA class
  \KOMAoptions{parskip=half}}
\makeatother
% Make \paragraph and \subparagraph free-standing
\makeatletter
\ifx\paragraph\undefined\else
  \let\oldparagraph\paragraph
  \renewcommand{\paragraph}{
    \@ifstar
      \xxxParagraphStar
      \xxxParagraphNoStar
  }
  \newcommand{\xxxParagraphStar}[1]{\oldparagraph*{#1}\mbox{}}
  \newcommand{\xxxParagraphNoStar}[1]{\oldparagraph{#1}\mbox{}}
\fi
\ifx\subparagraph\undefined\else
  \let\oldsubparagraph\subparagraph
  \renewcommand{\subparagraph}{
    \@ifstar
      \xxxSubParagraphStar
      \xxxSubParagraphNoStar
  }
  \newcommand{\xxxSubParagraphStar}[1]{\oldsubparagraph*{#1}\mbox{}}
  \newcommand{\xxxSubParagraphNoStar}[1]{\oldsubparagraph{#1}\mbox{}}
\fi
\makeatother

\usepackage{color}
\usepackage{fancyvrb}
\newcommand{\VerbBar}{|}
\newcommand{\VERB}{\Verb[commandchars=\\\{\}]}
\DefineVerbatimEnvironment{Highlighting}{Verbatim}{commandchars=\\\{\}}
% Add ',fontsize=\small' for more characters per line
\usepackage{framed}
\definecolor{shadecolor}{RGB}{241,243,245}
\newenvironment{Shaded}{\begin{snugshade}}{\end{snugshade}}
\newcommand{\AlertTok}[1]{\textcolor[rgb]{0.68,0.00,0.00}{#1}}
\newcommand{\AnnotationTok}[1]{\textcolor[rgb]{0.37,0.37,0.37}{#1}}
\newcommand{\AttributeTok}[1]{\textcolor[rgb]{0.40,0.45,0.13}{#1}}
\newcommand{\BaseNTok}[1]{\textcolor[rgb]{0.68,0.00,0.00}{#1}}
\newcommand{\BuiltInTok}[1]{\textcolor[rgb]{0.00,0.23,0.31}{#1}}
\newcommand{\CharTok}[1]{\textcolor[rgb]{0.13,0.47,0.30}{#1}}
\newcommand{\CommentTok}[1]{\textcolor[rgb]{0.37,0.37,0.37}{#1}}
\newcommand{\CommentVarTok}[1]{\textcolor[rgb]{0.37,0.37,0.37}{\textit{#1}}}
\newcommand{\ConstantTok}[1]{\textcolor[rgb]{0.56,0.35,0.01}{#1}}
\newcommand{\ControlFlowTok}[1]{\textcolor[rgb]{0.00,0.23,0.31}{\textbf{#1}}}
\newcommand{\DataTypeTok}[1]{\textcolor[rgb]{0.68,0.00,0.00}{#1}}
\newcommand{\DecValTok}[1]{\textcolor[rgb]{0.68,0.00,0.00}{#1}}
\newcommand{\DocumentationTok}[1]{\textcolor[rgb]{0.37,0.37,0.37}{\textit{#1}}}
\newcommand{\ErrorTok}[1]{\textcolor[rgb]{0.68,0.00,0.00}{#1}}
\newcommand{\ExtensionTok}[1]{\textcolor[rgb]{0.00,0.23,0.31}{#1}}
\newcommand{\FloatTok}[1]{\textcolor[rgb]{0.68,0.00,0.00}{#1}}
\newcommand{\FunctionTok}[1]{\textcolor[rgb]{0.28,0.35,0.67}{#1}}
\newcommand{\ImportTok}[1]{\textcolor[rgb]{0.00,0.46,0.62}{#1}}
\newcommand{\InformationTok}[1]{\textcolor[rgb]{0.37,0.37,0.37}{#1}}
\newcommand{\KeywordTok}[1]{\textcolor[rgb]{0.00,0.23,0.31}{\textbf{#1}}}
\newcommand{\NormalTok}[1]{\textcolor[rgb]{0.00,0.23,0.31}{#1}}
\newcommand{\OperatorTok}[1]{\textcolor[rgb]{0.37,0.37,0.37}{#1}}
\newcommand{\OtherTok}[1]{\textcolor[rgb]{0.00,0.23,0.31}{#1}}
\newcommand{\PreprocessorTok}[1]{\textcolor[rgb]{0.68,0.00,0.00}{#1}}
\newcommand{\RegionMarkerTok}[1]{\textcolor[rgb]{0.00,0.23,0.31}{#1}}
\newcommand{\SpecialCharTok}[1]{\textcolor[rgb]{0.37,0.37,0.37}{#1}}
\newcommand{\SpecialStringTok}[1]{\textcolor[rgb]{0.13,0.47,0.30}{#1}}
\newcommand{\StringTok}[1]{\textcolor[rgb]{0.13,0.47,0.30}{#1}}
\newcommand{\VariableTok}[1]{\textcolor[rgb]{0.07,0.07,0.07}{#1}}
\newcommand{\VerbatimStringTok}[1]{\textcolor[rgb]{0.13,0.47,0.30}{#1}}
\newcommand{\WarningTok}[1]{\textcolor[rgb]{0.37,0.37,0.37}{\textit{#1}}}

\usepackage{longtable,booktabs,array}
\newcounter{none} % for unnumbered tables
\usepackage{calc} % for calculating minipage widths
% Correct order of tables after \paragraph or \subparagraph
\usepackage{etoolbox}
\makeatletter
\patchcmd\longtable{\par}{\if@noskipsec\mbox{}\fi\par}{}{}
\makeatother
% Allow footnotes in longtable head/foot
\IfFileExists{footnotehyper.sty}{\usepackage{footnotehyper}}{\usepackage{footnote}}
\makesavenoteenv{longtable}
\usepackage{graphicx}
\makeatletter
\newsavebox\pandoc@box
\newcommand*\pandocbounded[1]{% scales image to fit in text height/width
  \sbox\pandoc@box{#1}%
  \Gscale@div\@tempa{\textheight}{\dimexpr\ht\pandoc@box+\dp\pandoc@box\relax}%
  \Gscale@div\@tempb{\linewidth}{\wd\pandoc@box}%
  \ifdim\@tempb\p@<\@tempa\p@\let\@tempa\@tempb\fi% select the smaller of both
  \ifdim\@tempa\p@<\p@\scalebox{\@tempa}{\usebox\pandoc@box}%
  \else\usebox{\pandoc@box}%
  \fi%
}
% Set default figure placement to htbp
\def\fps@figure{htbp}
\makeatother


% definitions for citeproc citations
\NewDocumentCommand\citeproctext{}{}
\NewDocumentCommand\citeproc{mm}{%
  \begingroup\def\citeproctext{#2}\cite{#1}\endgroup}
\makeatletter
 % allow citations to break across lines
 \let\@cite@ofmt\@firstofone
 % avoid brackets around text for \cite:
 \def\@biblabel#1{}
 \def\@cite#1#2{{#1\if@tempswa , #2\fi}}
\makeatother
\newlength{\cslhangindent}
\setlength{\cslhangindent}{1.5em}
\newlength{\csllabelwidth}
\setlength{\csllabelwidth}{3em}
\newenvironment{CSLReferences}[2] % #1 hanging-indent, #2 entry-spacing
 {\begin{list}{}{%
  \setlength{\itemindent}{0pt}
  \setlength{\leftmargin}{0pt}
  \setlength{\parsep}{0pt}
  % turn on hanging indent if param 1 is 1
  \ifodd #1
   \setlength{\leftmargin}{\cslhangindent}
   \setlength{\itemindent}{-1\cslhangindent}
  \fi
  % set entry spacing
  \setlength{\itemsep}{#2\baselineskip}}}
 {\end{list}}
\usepackage{calc}
\newcommand{\CSLBlock}[1]{\hfill\break\parbox[t]{\linewidth}{\strut\ignorespaces#1\strut}}
\newcommand{\CSLLeftMargin}[1]{\parbox[t]{\csllabelwidth}{\strut#1\strut}}
\newcommand{\CSLRightInline}[1]{\parbox[t]{\linewidth - \csllabelwidth}{\strut#1\strut}}
\newcommand{\CSLIndent}[1]{\hspace{\cslhangindent}#1}

\ifLuaTeX
\usepackage[bidi=basic,shorthands=off,provide=*]{babel}
\else
\usepackage[bidi=default,shorthands=off,provide=*]{babel}
\fi
\ifLuaTeX
  \usepackage{selnolig} % disable illegal ligatures
\fi


\setlength{\emergencystretch}{3em} % prevent overfull lines

\providecommand{\tightlist}{%
  \setlength{\itemsep}{0pt}\setlength{\parskip}{0pt}}



 


\KOMAoption{captions}{tableheading}
\makeatletter
\@ifpackageloaded{tcolorbox}{}{\usepackage[skins,breakable]{tcolorbox}}
\@ifpackageloaded{fontawesome5}{}{\usepackage{fontawesome5}}
\definecolor{quarto-callout-color}{HTML}{909090}
\definecolor{quarto-callout-note-color}{HTML}{0758E5}
\definecolor{quarto-callout-important-color}{HTML}{CC1914}
\definecolor{quarto-callout-warning-color}{HTML}{EB9113}
\definecolor{quarto-callout-tip-color}{HTML}{00A047}
\definecolor{quarto-callout-caution-color}{HTML}{FC5300}
\definecolor{quarto-callout-color-frame}{HTML}{acacac}
\definecolor{quarto-callout-note-color-frame}{HTML}{4582ec}
\definecolor{quarto-callout-important-color-frame}{HTML}{d9534f}
\definecolor{quarto-callout-warning-color-frame}{HTML}{f0ad4e}
\definecolor{quarto-callout-tip-color-frame}{HTML}{02b875}
\definecolor{quarto-callout-caution-color-frame}{HTML}{fd7e14}
\makeatother
\makeatletter
\@ifpackageloaded{bookmark}{}{\usepackage{bookmark}}
\makeatother
\makeatletter
\@ifpackageloaded{caption}{}{\usepackage{caption}}
\AtBeginDocument{%
\ifdefined\contentsname
  \renewcommand*\contentsname{목차}
\else
  \newcommand\contentsname{목차}
\fi
\ifdefined\listfigurename
  \renewcommand*\listfigurename{그림 목록}
\else
  \newcommand\listfigurename{그림 목록}
\fi
\ifdefined\listtablename
  \renewcommand*\listtablename{표 목록}
\else
  \newcommand\listtablename{표 목록}
\fi
\ifdefined\figurename
  \renewcommand*\figurename{그림}
\else
  \newcommand\figurename{그림}
\fi
\ifdefined\tablename
  \renewcommand*\tablename{표}
\else
  \newcommand\tablename{표}
\fi
}
\@ifpackageloaded{float}{}{\usepackage{float}}
\floatstyle{ruled}
\@ifundefined{c@chapter}{\newfloat{codelisting}{h}{lop}}{\newfloat{codelisting}{h}{lop}[chapter]}
\floatname{codelisting}{목록}
\newcommand*\listoflistings{\listof{codelisting}{코드 목록}}
\makeatother
\makeatletter
\makeatother
\makeatletter
\@ifpackageloaded{caption}{}{\usepackage{caption}}
\@ifpackageloaded{subcaption}{}{\usepackage{subcaption}}
\makeatother
\usepackage{bookmark}
\IfFileExists{xurl.sty}{\usepackage{xurl}}{} % add URL line breaks if available
\urlstyle{same}
% fallback for those not using the hyperref driver hyperxmp:
\makeatletter
\@ifundefined{xmpquote}{\newcommand{\xmpquote}[1]{#1}}{}
\makeatother
\hypersetup{
  pdftitle={컴퓨팅적사고와논리적문제해결},
  pdfauthor={신기홍},
  pdflang={ko},
  colorlinks=true,
  linkcolor={blue},
  filecolor={Maroon},
  citecolor={Blue},
  urlcolor={Blue},
  pdfcreator={LaTeX via pandoc}}


\title{컴퓨팅적사고와논리적문제해결}
\usepackage{etoolbox}
\makeatletter
\providecommand{\subtitle}[1]{% add subtitle to \maketitle
  \apptocmd{\@title}{\par {\large #1 \par}}{}{}
}
\makeatother
\subtitle{Computational Thinking \& Logical Problem Solving}
\author{신기홍}
\date{2026-01-01}
\begin{document}
\maketitle

\renewcommand*\contentsname{목차}
{
\setcounter{tocdepth}{2}
\tableofcontents
}

\bookmarksetup{startatroot}

\chapter{컴퓨팅적사고와논리적문제해결}\label{uxcef4uxd4e8uxd305uxc801uxc0acuxace0uxc640uxb17cuxb9acuxc801uxbb38uxc81cuxd574uxacb0}

이 강의는 \textbf{컴퓨팅 사고(Computational Thinking)}를 익히고
파이썬으로 직접 구현하는 능력을 기르는 것을 목표로 합니다.

\section{강의 개요}\label{uxac15uxc758-uxac1cuxc694}

{\def\LTcaptype{none} % do not increment counter
\begin{longtable}[]{@{}ll@{}}
\toprule\noalign{}
항목 & 내용 \\
\midrule\noalign{}
\endhead
\bottomrule\noalign{}
\endlastfoot
교과목명 & 컴퓨팅적사고와논리적문제해결 \\
담당교수 & 신기홍 \\
학점/시수 & 3학점 (이론 1 / 실습 2) \\
주교재 & 난생처음 컴퓨팅 사고 with 파이썬 (김명호, 한빛아카데미) \\
부교재 & AI 시대의 파이썬 데이터 분석 (오세종, 한빛아카데미) \\
\end{longtable}
}

\section{주차별 구성}\label{uxc8fcuxcc28uxbcc4-uxad6cuxc131}

각 주차는 \textbf{강의 본문 + 코딩 실습 + 연습문제} 세트로 구성됩니다.

\bookmarksetup{startatroot}

\chapter{컴퓨팅 사고와 디지털 문제해결}\label{sec-ch1}

\begin{quote}
이 장에서는 "컴퓨팅 사고"가 무엇이며, 왜 프로그래밍을 배우기 전에 먼저
익혀야 하는지 이해합니다. 문제 분해·패턴 인식·추상화·알고리즘이라는
{컴퓨팅 사고의 4단계}를 실생활 예시로 살펴보고, 컴퓨터와 소프트웨어의
역할, 소프트웨어 개발 6단계를 함께 파악합니다. 이 장에서 익히는 사고
습관은 앞으로 파이썬 코드를 작성하는 모든 순간의 뼈대가 됩니다.
\end{quote}

\section{왜 '코딩'이 아니라 '컴퓨팅 사고'인가}\label{sec-why-ct}

전 세계 많은 나라가 초·중·고등학교는 물론 대학 교양 과정에도
프로그래밍을 필수 과목으로 편성하고 있습니다. 그런데 그 목적은 단순히
``코드를 짤 줄 안다''는 기술 습득에 있지 않습니다. 진짜 목적은 {컴퓨팅
사고(Computational Thinking)}를 키우는 것입니다.

컴퓨팅 사고는 카네기멜론대학교의 컴퓨터과학자 Jeannette Wing이 2006년
\emph{Communications of the ACM}에 발표한 논문에서 체계화한
개념입니다(Wing 2006년). 그녀는 컴퓨팅 사고를 \textbf{``문제를
컴퓨터(또는 사람)가 효과적으로 실행할 수 있는 방식으로 공식화하고 그
해결책을 표현하는 사고 과정''}으로 정의하며, 읽기·쓰기·셈하기처럼 모든
현대인이 갖춰야 할 기본 소양이라고 주장했습니다.

\begin{tcolorbox}[enhanced jigsaw, arc=.35mm, bottomrule=.15mm, bottomtitle=1mm, breakable, colback=white, colbacktitle=quarto-callout-note-color!10!white, colframe=quarto-callout-note-color-frame, coltitle=black, left=2mm, leftrule=.75mm, opacityback=0, opacitybacktitle=0.6, rightrule=.15mm, title=\textcolor{quarto-callout-note-color}{\faInfo}\hspace{0.5em}{컴퓨팅 사고 ≠ 프로그래밍}, titlerule=0mm, toprule=.15mm, toptitle=1mm]

컴퓨팅 사고는 프로그래밍과 동의어가 아닙니다. 프로그래밍 없이도 컴퓨팅
사고를 연습할 수 있고, 반대로 프로그래밍을 해도 컴퓨팅 사고를 적용하지
않으면 좋은 코드를 만들기 어렵습니다. 컴퓨팅 사고는 \textbf{어떻게 풀
것인가}를 정리하는 과정이고, 프로그래밍은 그 생각을 컴퓨터에게 전달하는
수단입니다.

\end{tcolorbox}

Grover와 Pea는 K-12 컴퓨팅 사고 교육 연구에서 컴퓨팅 사고가
수학·과학·인문학 등 모든 분야에 걸쳐 전이될 수 있는 보편적 문제해결
능력임을 강조했습니다(Grover 와/과 Pea 2013년). CSTA(Computer Science
Teachers Association)의 K-12 컴퓨터과학 표준 역시 컴퓨팅 사고를 독립적인
핵심 역량으로 규정하고 있습니다(CSTA 2017년).

\begin{Shaded}
\begin{Highlighting}[]
\NormalTok{flowchart LR}
\NormalTok{    P["현실의 문제"] {-}{-}\textgreater{} CT["컴퓨팅 사고\textbackslash{}n(사고 과정)"]}
\NormalTok{    CT {-}{-}\textgreater{} PG["프로그래밍\textbackslash{}n(표현 수단)"]}
\NormalTok{    PG {-}{-}\textgreater{} S["컴퓨터의\textbackslash{}n실행과 해결"]}
\NormalTok{    style CT fill:\#dbeafe}
\NormalTok{    style PG fill:\#fde68a}
\end{Highlighting}
\end{Shaded}

\section{컴퓨터와 소프트웨어의 역할}\label{sec-hw-sw}

컴퓨팅 사고를 이해하기 전에, 우리가 다루는 도구인 컴퓨터의 구조를 간단히
살펴봅니다.

\subsection{하드웨어와
소프트웨어}\label{uxd558uxb4dcuxc6e8uxc5b4uxc640-uxc18cuxd504uxd2b8uxc6e8uxc5b4}

컴퓨터는 크게 \textbf{하드웨어(Hardware)}와
\textbf{소프트웨어(Software)}로 이루어집니다.

\begin{itemize}
\tightlist
\item
  \textbf{하드웨어}: 컴퓨터를 구성하는 물리적 장치입니다.
  중앙처리장치(CPU), 기억장치(메모리·저장소), 입력장치(키보드·마우스),
  출력장치(모니터·프린터) 등이 여기 해당합니다.
\item
  \textbf{소프트웨어}: 하드웨어를 작동시키는 명령어의 집합으로, 보통
  \textbf{프로그램(Program)}이라고도 부릅니다. 소프트웨어는 먼저
  메모리에 탑재된 다음, 기계어로 변환되어 CPU와 각종 장치를
  작동시킵니다.
\end{itemize}

\begin{Shaded}
\begin{Highlighting}[]
\NormalTok{flowchart TD}
\NormalTok{    SW["소프트웨어\textbackslash{}n(프로그램)"] {-}{-}\textgreater{}|"메모리 탑재"| MEM["메모리"]}
\NormalTok{    MEM {-}{-}\textgreater{}|"기계어 변환"| CPU["중앙처리장치\textbackslash{}n(CPU)"]}
\NormalTok{    CPU {-}{-}\textgreater{} IO["입출력 장치\textbackslash{}n(키보드 / 모니터 등)"]}
\NormalTok{    style SW fill:\#dbeafe}
\NormalTok{    style CPU fill:\#fde68a}
\end{Highlighting}
\end{Shaded}

\subsection{프로그래밍
언어}\label{uxd504uxb85cuxadf8uxb798uxbc0d-uxc5b8uxc5b4}

컴퓨터는 0과 1로 이루어진 \textbf{기계어(machine code)}만 이해합니다.
하지만 사람이 직접 0과 1을 나열하는 것은 매우 어렵기 때문에, 인간의
언어와 유사한 \textbf{고급 언어(high-level language)}가 발명되었습니다.

{\def\LTcaptype{none} % do not increment counter
\begin{longtable}[]{@{}
  >{\raggedright\arraybackslash}p{(\linewidth - 4\tabcolsep) * \real{0.3333}}
  >{\raggedright\arraybackslash}p{(\linewidth - 4\tabcolsep) * \real{0.3333}}
  >{\raggedright\arraybackslash}p{(\linewidth - 4\tabcolsep) * \real{0.3333}}@{}}
\toprule\noalign{}
\begin{minipage}[b]{\linewidth}\raggedright
분류
\end{minipage} & \begin{minipage}[b]{\linewidth}\raggedright
특징
\end{minipage} & \begin{minipage}[b]{\linewidth}\raggedright
예시
\end{minipage} \\
\midrule\noalign{}
\endhead
\bottomrule\noalign{}
\endlastfoot
저급 언어 & 기계에 가까워 빠르지만 배우기 어렵고 유지보수 힘듦 & 기계어,
어셈블리어 \\
고급 언어 & 사람 언어와 유사, 배우기 쉽고 유지보수 용이 & C, Java,
\textbf{Python} 등 \\
\end{longtable}
}

고급 언어로 작성한 코드를 \textbf{소스코드(source code)}라 하고, 이를
컴퓨터가 실행할 수 있는 형태로 변환하는 도구를
\textbf{컴파일러(compiler)} 또는 \textbf{인터프리터(interpreter)}라
합니다. 파이썬은 대표적인 \textbf{인터프리터 방식}의 고급 언어로, 한
줄씩 즉시 해석하고 실행합니다.

\begin{tcolorbox}[enhanced jigsaw, arc=.35mm, bottomrule=.15mm, bottomtitle=1mm, breakable, colback=white, colbacktitle=quarto-callout-tip-color!10!white, colframe=quarto-callout-tip-color-frame, coltitle=black, left=2mm, leftrule=.75mm, opacityback=0, opacitybacktitle=0.6, rightrule=.15mm, title=\textcolor{quarto-callout-tip-color}{\faLightbulb}\hspace{0.5em}{왜 파이썬인가?}, titlerule=0mm, toprule=.15mm, toptitle=1mm]

파이썬은 문법이 영어 문장과 유사해 배우기 쉽고, 데이터 분석·인공지능·웹
개발 등 다양한 분야에서 가장 널리 사용되는 언어입니다. Stack Overflow
개발자 설문(2023)에서도 10년 이상 연속으로 ``가장 많이 배우고 싶은
언어'' 상위권을 유지하고 있습니다.

\end{tcolorbox}

\section{소프트웨어 개발 6단계}\label{sec-sw-dev}

실제 소프트웨어는 생각나는 대로 코드를 쓰는 것이 아니라, 정해진
\textbf{개발 과정(process)}을 따릅니다. 규모가 클수록 많은 인원이
투입되므로, 체계적인 프로세스가 필수입니다.

\begin{Shaded}
\begin{Highlighting}[]
\NormalTok{flowchart LR}
\NormalTok{    A["1단계\textbackslash{}n문제 분석"] {-}{-}\textgreater{} B["2단계\textbackslash{}n설계"]}
\NormalTok{    B {-}{-}\textgreater{} C["3단계\textbackslash{}n프로그래밍"]}
\NormalTok{    C {-}{-}\textgreater{} D["4단계\textbackslash{}n테스트"]}
\NormalTok{    D {-}{-}\textgreater{} E["5단계\textbackslash{}n사용"]}
\NormalTok{    E {-}{-}\textgreater{} F["6단계\textbackslash{}n유지보수"]}
\NormalTok{    F {-}.{-}\textgreater{}|"개선 요구"| A}
\NormalTok{    style A fill:\#dbeafe}
\NormalTok{    style C fill:\#fde68a}
\NormalTok{    style F fill:\#dcfce7}
\end{Highlighting}
\end{Shaded}

{\def\LTcaptype{none} % do not increment counter
\begin{longtable}[]{@{}clll@{}}
\toprule\noalign{}
단계 & 이름 & 핵심 질문 & 주요 활동 \\
\midrule\noalign{}
\endhead
\bottomrule\noalign{}
\endlastfoot
1 & 문제 분석 & \textbf{무엇을} 만들 것인가? & 요구사항 수집, 명세서
작성 \\
2 & 설계 & \textbf{어떻게} 만들 것인가? & 알고리즘 설계, 구조 결정 \\
3 & 프로그래밍 & 코드로 \textbf{구현} & 파이썬 등 언어로 작성 \\
4 & 테스트 & \textbf{제대로 동작}하는가? & 다양한 환경에서 검증 \\
5 & 사용 & 사용자에게 \textbf{배포} & 매뉴얼 제공, 교육 \\
6 & 유지보수 & \textbf{개선·업그레이드} & 버그 수정, 기능 추가 \\
\end{longtable}
}

\begin{tcolorbox}[enhanced jigsaw, arc=.35mm, bottomrule=.15mm, bottomtitle=1mm, breakable, colback=white, colbacktitle=quarto-callout-important-color!10!white, colframe=quarto-callout-important-color-frame, coltitle=black, left=2mm, leftrule=.75mm, opacityback=0, opacitybacktitle=0.6, rightrule=.15mm, title=\textcolor{quarto-callout-important-color}{\faExclamation}\hspace{0.5em}{2단계 설계 = 컴퓨팅 사고의 실천}, titlerule=0mm, toprule=.15mm, toptitle=1mm]

소프트웨어 개발 6단계에서 \textbf{설계(2단계)}가 바로 컴퓨팅 사고가 가장
직접적으로 쓰이는 단계입니다. 문제를 어떻게 쪼개고, 어떤 순서로 해결할지
알고리즘으로 구체화하는 과정 전체가 컴퓨팅 사고의 실천입니다.

\end{tcolorbox}

\section{컴퓨팅 사고의 4단계}\label{sec-ct-4steps}

컴퓨팅 사고는 복잡한 문제를 컴퓨터가 처리할 수 있는 방식으로 풀어가는 네
가지 사고 단계로 구성됩니다(Wing 2006년; Grover 와/과 Pea 2013년).

\begin{Shaded}
\begin{Highlighting}[]
\NormalTok{flowchart LR}
\NormalTok{    D["1 문제 분해\textbackslash{}nDecomposition"]}
\NormalTok{    P["2 패턴 인식\textbackslash{}nPattern Recognition"]}
\NormalTok{    A["3 추상화\textbackslash{}nAbstraction"]}
\NormalTok{    G["4 알고리즘\textbackslash{}nAlgorithm"]}
\NormalTok{    D {-}{-}\textgreater{} P {-}{-}\textgreater{} A {-}{-}\textgreater{} G}
\NormalTok{    G {-}.{-}\textgreater{}|"새 문제에 적용"| D}
\end{Highlighting}
\end{Shaded}

\subsection{1단계: 문제 분해
(Decomposition)}\label{uxb2e8uxacc4-uxbb38uxc81c-uxbd84uxd574-decomposition}

{문제 분해}란 크고 복잡한 문제를 \textbf{작고 관리하기 쉬운 여러 개의
부분 문제}로 나누는 것입니다. 큰 문제를 한 번에 해결하려 하면
막막하지만, 작은 조각으로 나누면 각각을 순서대로 또는 동시에 해결할 수
있습니다.

\begin{Shaded}
\begin{Highlighting}[]
\NormalTok{flowchart TD}
\NormalTok{    PROB["저녁 식사 준비"] {-}{-}\textgreater{} A["재료 목록 작성"]}
\NormalTok{    PROB {-}{-}\textgreater{} B["마트 장보기"]}
\NormalTok{    PROB {-}{-}\textgreater{} C["요리하기"]}
\NormalTok{    B {-}{-}\textgreater{} B1["채소 구매"]}
\NormalTok{    B {-}{-}\textgreater{} B2["육류 구매"]}
\NormalTok{    B {-}{-}\textgreater{} B3["양념 구매"]}
\NormalTok{    C {-}{-}\textgreater{} C1["재료 손질"]}
\NormalTok{    C {-}{-}\textgreater{} C2["조리"]}
\NormalTok{    C {-}{-}\textgreater{} C3["상차림"]}
\NormalTok{    style PROB fill:\#fde68a}
\NormalTok{    style A fill:\#dbeafe}
\NormalTok{    style B fill:\#dbeafe}
\NormalTok{    style C fill:\#dbeafe}
\end{Highlighting}
\end{Shaded}

\textbf{그림 1.1} 저녁 식사 준비 문제의 분해 예시. 하나의 큰 문제가 여러
작은 문제로 나뉘고, 일부는 순차적으로, 일부(장보기 항목들)는 병렬적으로
해결됩니다.

\textbf{핵심.} 문제 분해의 결과물은 각각 독립적으로 해결 가능한 작은
문제들입니다. 작은 문제의 해답을 합치면 큰 문제가 해결됩니다.

\subsection{2단계: 패턴 인식 (Pattern
Recognition)}\label{uxb2e8uxacc4-uxd328uxd134-uxc778uxc2dd-pattern-recognition}

{패턴 인식}이란 분해된 여러 문제 또는 데이터 속에서 \textbf{반복되는
공통점, 규칙, 추세}를 찾아내는 단계입니다. 공통된 패턴을 발견하면 같은
해법을 여러 문제에 재사용할 수 있어 효율이 크게 높아집니다.

\textbf{실생활 예시:} 마트에서 물건을 여러 개 살 때, 각 물건의 가격을
더하는 과정은 항상 동일합니다. ``(수량) × (단가)'' 패턴을 발견하면 물건
종류가 달라져도 같은 계산 방법을 적용할 수 있습니다.

\textbf{프로그래밍에서의 패턴 인식:}

\begin{itemize}
\tightlist
\item
  비슷한 작업이 여러 번 반복될 때 → \textbf{반복문(loop)}으로 처리
\item
  조건에 따라 다른 행동이 필요할 때 → \textbf{조건문(if)}으로 처리
\item
  자주 쓰는 코드 묶음 → \textbf{함수(function)}로 재사용
\end{itemize}

\subsection{3단계: 추상화
(Abstraction)}\label{uxb2e8uxacc4-uxcd94uxc0c1uxd654-abstraction}

{추상화}란 문제를 해결하는 데 \textbf{꼭 필요한 정보만 남기고, 중요하지
않은 세부 사항은 무시}하는 것입니다. 복잡한 현실에서 핵심만 뽑아
단순하고 일반적인 모델을 만드는 과정입니다(Wing 2006년).

패턴 인식을 통해 발견한 해법을 \textbf{일반화(generalization)}하면,
유사한 상황에 폭넓게 적용할 수 있는 모델이 됩니다.

\begin{quote}
\textbf{구체적 상황}: 사과 3개(개당 500원), 바나나 2개(개당 800원), 우유
1개(개당 2000원)의 총 가격은?

\textbf{추상화된 공식}: 총 가격 = Σ (수량 × 단가)

→ 물건 종류와 개수가 바뀌어도 이 공식은 항상 성립합니다.
\end{quote}

{\def\LTcaptype{none} % do not increment counter
\begin{longtable}[]{@{}
  >{\raggedleft\arraybackslash}p{(\linewidth - 2\tabcolsep) * \real{0.5000}}
  >{\raggedright\arraybackslash}p{(\linewidth - 2\tabcolsep) * \real{0.5000}}@{}}
\toprule\noalign{}
\begin{minipage}[b]{\linewidth}\raggedleft
추상화 전 (구체)
\end{minipage} & \begin{minipage}[b]{\linewidth}\raggedright
추상화 후 (일반화)
\end{minipage} \\
\midrule\noalign{}
\endhead
\bottomrule\noalign{}
\endlastfoot
사과 3 × 500 + 바나나 2 × 800 + \ldots{} & 총합 = Σ (수량ᵢ × 단가ᵢ) \\
등기는 등기 창구, 택배는 택배 창구 & 우편물은 종류에 맞는 창구에 접수 \\
반복문으로 1\textasciitilde100 합산 & \texttt{sum(range(1,\ 101))} \\
\end{longtable}
}

\begin{tcolorbox}[enhanced jigsaw, arc=.35mm, bottomrule=.15mm, bottomtitle=1mm, breakable, colback=white, colbacktitle=quarto-callout-note-color!10!white, colframe=quarto-callout-note-color-frame, coltitle=black, left=2mm, leftrule=.75mm, opacityback=0, opacitybacktitle=0.6, rightrule=.15mm, title=\textcolor{quarto-callout-note-color}{\faInfo}\hspace{0.5em}{추상화 = 프로그래밍의 핵심 기술}, titlerule=0mm, toprule=.15mm, toptitle=1mm]

파이썬에서 함수(function)·클래스(class)·모듈(module)은 모두 추상화의
도구입니다. 복잡한 처리를 하나의 이름으로 묶어 필요할 때마다 꺼내 쓸 수
있게 합니다.

\end{tcolorbox}

\subsection{4단계: 알고리즘
(Algorithm)}\label{uxb2e8uxacc4-uxc54cuxace0uxb9acuxc998-algorithm}

{알고리즘}이란 특정 문제를 해결하기 위한 \textbf{명확하고 순서가 있는
명령어의 집합}입니다. 앞의 세 단계를 통해 문제를 정리했다면, 마지막으로
컴퓨터가 실행할 수 있는 구체적인 절차를 설계하는 것이 알고리즘입니다.

좋은 알고리즘의 세 가지 요건:

\begin{enumerate}
\def\labelenumi{\arabic{enumi}.}
\tightlist
\item
  \textbf{명확성}: 각 단계가 모호하지 않아야 합니다.
\item
  \textbf{유한성}: 반드시 끝나야 합니다 (무한 반복 금지).
\item
  \textbf{효율성}: 같은 결과라면 더 빠르고 자원을 덜 쓰는 것이 좋습니다.
\end{enumerate}

\begin{Shaded}
\begin{Highlighting}[]
\NormalTok{flowchart TD}
\NormalTok{    START(["시작"]) {-}{-}\textgreater{} A["냄비에 물 550ml 붓기"]}
\NormalTok{    A {-}{-}\textgreater{} B["물 끓이기"]}
\NormalTok{    B {-}{-}\textgreater{} C\{"물이 끓었는가?"\}}
\NormalTok{    C {-}{-}\textgreater{}|"아니오"| B}
\NormalTok{    C {-}{-}\textgreater{}|"예"| D["면 · 스프 · 후레이크 넣기"]}
\NormalTok{    D {-}{-}\textgreater{} E["4분 30초 더 끓이기"]}
\NormalTok{    E {-}{-}\textgreater{} F["그릇에 담기"]}
\NormalTok{    F {-}{-}\textgreater{} END(["완료"])}
\NormalTok{    style START fill:\#dcfce7}
\NormalTok{    style END fill:\#fde68a}
\NormalTok{    style C fill:\#fde68a}
\end{Highlighting}
\end{Shaded}

\textbf{그림 1.2} ``라면 끓이기''를 순서도(flowchart)로 표현한 알고리즘.
조건 판단(마름모)과 반복(화살표 되돌아오기)이 포함되어 있습니다.

\section{컴퓨팅 사고 4단계의 통합: 실전 예제}\label{sec-ct-example}

지금까지 배운 4단계를 하나의 실전 문제에 적용해 봅니다.

\textbf{문제}: 반 학생 30명의 수학 시험 점수가 주어졌을 때,
\textbf{평균보다 높은 점수}를 받은 학생이 몇 명인지 구하시오.

\subsection{1단계 적용: 문제
분해}\label{uxb2e8uxacc4-uxc801uxc6a9-uxbb38uxc81c-uxbd84uxd574}

\begin{enumerate}
\def\labelenumi{\arabic{enumi}.}
\tightlist
\item
  점수 데이터를 입력받기
\item
  전체 평균을 계산하기
\item
  각 점수와 평균을 비교하기
\item
  평균보다 높은 학생 수를 세기
\item
  결과를 출력하기
\end{enumerate}

\subsection{2단계 적용: 패턴
인식}\label{uxb2e8uxacc4-uxc801uxc6a9-uxd328uxd134-uxc778uxc2dd}

\begin{itemize}
\tightlist
\item
  ``각 점수와 평균을 비교하기''는 30번 반복됩니다 → \textbf{반복 패턴} →
  반복문 사용
\item
  비교 결과에 따라 카운트 여부가 달라집니다 → \textbf{조건 패턴} →
  조건문 사용
\end{itemize}

\subsection{3단계 적용:
추상화}\label{uxb2e8uxacc4-uxc801uxc6a9-uxcd94uxc0c1uxd654}

\begin{itemize}
\tightlist
\item
  30명이라는 구체적인 숫자 대신 \textbf{``n명의 점수 목록''}으로 일반화
\item
  수학 점수뿐만 아니라 어떤 과목의 점수에도 같은 방법 적용 가능
\end{itemize}

\subsection{4단계 적용: 알고리즘 설계
(의사코드)}\label{uxb2e8uxacc4-uxc801uxc6a9-uxc54cuxace0uxb9acuxc998-uxc124uxacc4-uxc758uxc0acuxcf54uxb4dc}

\begin{verbatim}
입력: 점수 목록 scores
1. total ← 0
2. scores의 각 점수(s)에 대해:
     total ← total + s
3. avg ← total / 점수 개수
4. count ← 0
5. scores의 각 점수(s)에 대해:
     만약 s > avg 이면:
         count ← count + 1
6. count 출력
\end{verbatim}

이 알고리즘을 파이썬 코드로 옮기면 다음과 같습니다. 아직 문법을 몰라도
괜찮습니다. 앞으로 4주에 걸쳐 이 코드의 각 부분을 하나씩 배웁니다.

\begin{Shaded}
\begin{Highlighting}[]
\CommentTok{\# 예시 데이터 (30명의 점수)}
\NormalTok{scores }\OperatorTok{=}\NormalTok{ [}\DecValTok{78}\NormalTok{, }\DecValTok{91}\NormalTok{, }\DecValTok{65}\NormalTok{, }\DecValTok{82}\NormalTok{, }\DecValTok{75}\NormalTok{, }\DecValTok{88}\NormalTok{, }\DecValTok{70}\NormalTok{, }\DecValTok{95}\NormalTok{, }\DecValTok{60}\NormalTok{, }\DecValTok{83}\NormalTok{,}
          \DecValTok{77}\NormalTok{, }\DecValTok{89}\NormalTok{, }\DecValTok{72}\NormalTok{, }\DecValTok{68}\NormalTok{, }\DecValTok{91}\NormalTok{, }\DecValTok{84}\NormalTok{, }\DecValTok{76}\NormalTok{, }\DecValTok{80}\NormalTok{, }\DecValTok{55}\NormalTok{, }\DecValTok{93}\NormalTok{,}
          \DecValTok{74}\NormalTok{, }\DecValTok{87}\NormalTok{, }\DecValTok{69}\NormalTok{, }\DecValTok{85}\NormalTok{, }\DecValTok{78}\NormalTok{, }\DecValTok{92}\NormalTok{, }\DecValTok{63}\NormalTok{, }\DecValTok{79}\NormalTok{, }\DecValTok{88}\NormalTok{, }\DecValTok{71}\NormalTok{]}

\CommentTok{\# 평균 계산}
\NormalTok{avg }\OperatorTok{=} \BuiltInTok{sum}\NormalTok{(scores) }\OperatorTok{/} \BuiltInTok{len}\NormalTok{(scores)}
\BuiltInTok{print}\NormalTok{(}\SpecialStringTok{f"평균 점수: }\SpecialCharTok{\{}\NormalTok{avg}\SpecialCharTok{:.2f\}}\SpecialStringTok{"}\NormalTok{)}

\CommentTok{\# 평균보다 높은 학생 수 계산}
\NormalTok{count }\OperatorTok{=} \DecValTok{0}
\ControlFlowTok{for}\NormalTok{ s }\KeywordTok{in}\NormalTok{ scores:}
    \ControlFlowTok{if}\NormalTok{ s }\OperatorTok{\textgreater{}}\NormalTok{ avg:}
\NormalTok{        count }\OperatorTok{+=} \DecValTok{1}

\BuiltInTok{print}\NormalTok{(}\SpecialStringTok{f"평균 초과 학생 수: }\SpecialCharTok{\{}\NormalTok{count}\SpecialCharTok{\}}\SpecialStringTok{명"}\NormalTok{)}
\end{Highlighting}
\end{Shaded}

\begin{tcolorbox}[enhanced jigsaw, arc=.35mm, bottomrule=.15mm, bottomtitle=1mm, breakable, colback=white, colbacktitle=quarto-callout-tip-color!10!white, colframe=quarto-callout-tip-color-frame, coltitle=black, left=2mm, leftrule=.75mm, opacityback=0, opacitybacktitle=0.6, rightrule=.15mm, title=\textcolor{quarto-callout-tip-color}{\faLightbulb}\hspace{0.5em}{코드를 보는 관점}, titlerule=0mm, toprule=.15mm, toptitle=1mm]

위 코드에서 \texttt{for\ s\ in\ scores:}는 \textbf{반복 패턴}을,
\texttt{if\ s\ \textgreater{}\ avg:}는 \textbf{조건 패턴}을 구현한
것입니다. 코드를 볼 때 ``이 부분이 컴퓨팅 사고의 어느 단계에
해당하는가?''를 생각하는 습관을 들이면 프로그래밍 실력이 빠르게
향상됩니다.

\end{tcolorbox}

\section{순서도로 알고리즘 표현하기}\label{sec-flowchart}

알고리즘을 시각적으로 표현할 때 가장 많이 쓰이는 도구가
{순서도(flowchart)}입니다. 순서도는 표준화된 도형을 사용해 프로그램의
흐름을 누구나 이해할 수 있게 표현합니다.

{\def\LTcaptype{none} % do not increment counter
\begin{longtable}[]{@{}lll@{}}
\toprule\noalign{}
도형 & 의미 & 예시 \\
\midrule\noalign{}
\endhead
\bottomrule\noalign{}
\endlastfoot
타원 & 시작 / 끝 & 시작, 완료 \\
직사각형 & 처리 (명령 실행) & \texttt{count\ +=\ 1} \\
마름모 & 조건 판단 (예/아니오) & \texttt{s\ \textgreater{}\ avg?} \\
평행사변형 & 입력 / 출력 & 값 입력, 결과 출력 \\
화살표 & 흐름의 방향 & → \\
\end{longtable}
}

\begin{Shaded}
\begin{Highlighting}[]
\NormalTok{flowchart TD}
\NormalTok{    START(["시작"]) {-}{-}\textgreater{} IN[/"점수 목록 scores 입력"/]}
\NormalTok{    IN {-}{-}\textgreater{} AVG["avg = sum(scores) / len(scores)"]}
\NormalTok{    AVG {-}{-}\textgreater{} INIT["count = 0 · i = 0"]}
\NormalTok{    INIT {-}{-}\textgreater{} LOOP\{"i \textless{} len(scores)?"\}}
\NormalTok{    LOOP {-}{-}\textgreater{}|"아니오"| OUT[/"count 출력"/]}
\NormalTok{    LOOP {-}{-}\textgreater{}|"예"| COND\{"scores[i] \textgreater{} avg?"\}}
\NormalTok{    COND {-}{-}\textgreater{}|"예"| INC["count += 1"]}
\NormalTok{    COND {-}{-}\textgreater{}|"아니오"| NEXT["i += 1"]}
\NormalTok{    INC {-}{-}\textgreater{} NEXT}
\NormalTok{    NEXT {-}{-}\textgreater{} LOOP}
\NormalTok{    OUT {-}{-}\textgreater{} END(["완료"])}
\NormalTok{    style START fill:\#dcfce7}
\NormalTok{    style END fill:\#dcfce7}
\NormalTok{    style LOOP fill:\#fde68a}
\NormalTok{    style COND fill:\#fde68a}
\end{Highlighting}
\end{Shaded}

\textbf{그림 1.3} ``평균 초과 학생 수'' 알고리즘의 순서도.
마름모(조건)와 반복 구조(화살표가 위로 돌아오는 부분)가 프로그램의 핵심
논리를 시각적으로 보여 줍니다.

\section{좋은 알고리즘이란}\label{sec-good-algorithm}

같은 문제를 해결하더라도 알고리즘의 \textbf{효율성}은 크게 다를 수
있습니다.

\begin{Shaded}
\begin{Highlighting}[]
\CommentTok{\# 방법 1: 반복문으로 하나씩 더하기 (100번의 덧셈)}
\NormalTok{total }\OperatorTok{=} \DecValTok{0}
\ControlFlowTok{for}\NormalTok{ i }\KeywordTok{in} \BuiltInTok{range}\NormalTok{(}\DecValTok{1}\NormalTok{, }\DecValTok{101}\NormalTok{):}
\NormalTok{    total }\OperatorTok{+=}\NormalTok{ i}
\BuiltInTok{print}\NormalTok{(total)   }\CommentTok{\# 5050}

\CommentTok{\# 방법 2: 수학 공식 사용 (단 1번의 계산)}
\NormalTok{n }\OperatorTok{=} \DecValTok{100}
\NormalTok{total }\OperatorTok{=}\NormalTok{ n }\OperatorTok{*}\NormalTok{ (n }\OperatorTok{+} \DecValTok{1}\NormalTok{) }\OperatorTok{//} \DecValTok{2}
\BuiltInTok{print}\NormalTok{(total)   }\CommentTok{\# 5050}
\end{Highlighting}
\end{Shaded}

두 방법 모두 정답(5050)을 내지만, 방법 2는 숫자가 아무리 커져도 단 한
번의 계산으로 끝납니다. 이처럼 문제를 수학적으로 추상화하면 훨씬
효율적인 알고리즘을 설계할 수 있습니다.

\textbf{핵심.} 컴퓨팅 사고의 4단계는 순서대로 한 번만 적용하는 것이
아니라, 문제를 만날 때마다 반복적·순환적으로 활용합니다. 실생활의 어떤
문제도 \textbf{분해 → 패턴 인식 → 추상화 → 알고리즘} 과정을 거쳐
컴퓨터가 이해할 수 있는 형태로 바꿀 수 있습니다.

\section{이 강의에서 배울 것들}\label{sec-roadmap}

이 강의는 컴퓨팅 사고를 익히고, 파이썬으로 직접 구현하는 능력을 기르는
것을 목표로 합니다. 전체 15주의 학습 흐름은 다음과 같습니다.

\begin{Shaded}
\begin{Highlighting}[]
\NormalTok{flowchart TD}
\NormalTok{    subgraph A["전반부 (1–8주): 컴퓨팅 사고 \& 파이썬 기초"]}
\NormalTok{        A1["CT 개념 (1주)"] {-}{-}\textgreater{} A2["문제 분해 (2주)"] {-}{-}\textgreater{} A3["추상화·알고리즘 (3주)"]}
\NormalTok{        A3 {-}{-}\textgreater{} A4["데이터·변수 (4주)"] {-}{-}\textgreater{} A5["조건문 (5주)"]}
\NormalTok{        A5 {-}{-}\textgreater{} A6["반복문 (6주)"] {-}{-}\textgreater{} A7["자료구조 (7주)"] {-}{-}\textgreater{} A8["중간고사 (8주)"]}
\NormalTok{    end}
\NormalTok{    subgraph B["후반부 (9–15주): 데이터 분석 \& AI"]}
\NormalTok{        B1["함수 (9주)"] {-}{-}\textgreater{} B2["데이터 분석 입문 (10주)"] {-}{-}\textgreater{} B3["시각화 (11주)"]}
\NormalTok{        B3 {-}{-}\textgreater{} B4["패턴·관계 (12주)"] {-}{-}\textgreater{} B5["AI 원리 (13주)"]}
\NormalTok{        B5 {-}{-}\textgreater{} B6["생성형 AI (14주)"] {-}{-}\textgreater{} B7["프로젝트 (15주)"]}
\NormalTok{    end}
\NormalTok{    A {-}{-}\textgreater{} B}
\NormalTok{    style A fill:\#dbeafe,stroke:\#93c5fd}
\NormalTok{    style B fill:\#fef9c3,stroke:\#fbbf24}
\end{Highlighting}
\end{Shaded}

\textbf{그림 1.4} 15주 강의 로드맵. 전반부는 파이썬 기초 문법과 컴퓨팅
사고를, 후반부는 데이터 분석과 AI 활용을 다룹니다.

\begin{center}\rule{0.5\linewidth}{0.5pt}\end{center}

\section*{참고문헌}\label{uxcc38uxace0uxbb38uxd5cc}
\addcontentsline{toc}{section}{참고문헌}

\markright{참고문헌}

\protect\phantomsection\label{refs}
\begin{CSLReferences}{1}{1}
\bibitem[\citeproctext]{ref-csta2017standards}
CSTA. 2017년. \emph{CSTA K--12 Computer Science Standards}. Computer
Science Teachers Association.
\url{https://www.csteachers.org/page/standards}.

\bibitem[\citeproctext]{ref-grover2013computational}
Grover, Shuchi, 와/과 Roy Pea. 2013년. {``Computational Thinking in
K--12: A Review of the State of the Field''}. \emph{Educational
Researcher} 42 (1): 38--43.
\url{https://doi.org/10.3102/0013189X12463051}.

\bibitem[\citeproctext]{ref-wing2006computational}
Wing, Jeannette M. 2006년. {``Computational thinking''}.
\emph{Communications of the ACM} 49 (3): 33--35.
\url{https://doi.org/10.1145/1118178.1118215}.

\end{CSLReferences}

\bookmarksetup{startatroot}

\chapter{1주차 코딩 실습 (약 1시간)}\label{sec-ch1-lab}

이 실습은 1장에서 익힌 컴퓨팅 사고의 4단계(문제 분해·패턴
인식·추상화·알고리즘)를 파이썬 코드로 직접 확인하는 과정입니다.
\textbf{예제 코드}를 먼저 실행해 동작을 확인한 뒤, \textbf{제출용
문제}의 빈칸(\texttt{\#\ TODO})을 채워 제출합니다.

\begin{tcolorbox}[enhanced jigsaw, arc=.35mm, bottomrule=.15mm, bottomtitle=1mm, breakable, colback=white, colbacktitle=quarto-callout-note-color!10!white, colframe=quarto-callout-note-color-frame, coltitle=black, left=2mm, leftrule=.75mm, opacityback=0, opacitybacktitle=0.6, rightrule=.15mm, title=\textcolor{quarto-callout-note-color}{\faInfo}\hspace{0.5em}{실습 환경 확인}, titlerule=0mm, toprule=.15mm, toptitle=1mm]

파이썬이 설치되어 있지 않다면
\href{https://www.python.org/downloads/}{python.org}에서 설치합니다. VS
Code 또는 주피터 노트북(Jupyter Notebook) 중 편한 환경을 사용하세요.
\texttt{print("안녕,\ 파이썬!")} 한 줄을 실행해 화면에 출력이 나오면
준비 완료입니다.

\end{tcolorbox}

\section{예제 코드 (함께
실행)}\label{uxc608uxc81c-uxcf54uxb4dc-uxd568uxaed8-uxc2e4uxd589}

\subsection{예제 1: 컴퓨팅 사고 4단계로 ``최고 점수 학생
찾기''}\label{uxc608uxc81c-1-uxcef4uxd4e8uxd305-uxc0acuxace0-4uxb2e8uxacc4uxb85c-uxcd5cuxace0-uxc810uxc218-uxd559uxc0dd-uxcc3euxae30}

아래는 반 학생 5명의 이름과 점수가 주어졌을 때, CT 4단계를 적용해 최고
점수 학생을 찾는 완성된 예제입니다.

\begin{Shaded}
\begin{Highlighting}[]
\CommentTok{\# 데이터 정의}
\NormalTok{students }\OperatorTok{=}\NormalTok{ [}\StringTok{"김민준"}\NormalTok{, }\StringTok{"이서연"}\NormalTok{, }\StringTok{"박지훈"}\NormalTok{, }\StringTok{"최유나"}\NormalTok{, }\StringTok{"정도현"}\NormalTok{]}
\NormalTok{scores   }\OperatorTok{=}\NormalTok{ [}\DecValTok{85}\NormalTok{, }\DecValTok{92}\NormalTok{, }\DecValTok{78}\NormalTok{, }\DecValTok{96}\NormalTok{, }\DecValTok{88}\NormalTok{]}

\CommentTok{\# {-}{-}{-} 1단계: 문제 분해 {-}{-}{-}}
\CommentTok{\# (1) 최고 점수 찾기  (2) 해당 학생 이름 찾기  (3) 출력하기}

\CommentTok{\# {-}{-}{-} 2단계: 패턴 인식 {-}{-}{-}}
\CommentTok{\# 모든 점수를 한 번씩 확인해야 함 → 반복 패턴 → for 문 사용}
\CommentTok{\# 현재 최고보다 높으면 갱신 → 조건 패턴 → if 문 사용}

\CommentTok{\# {-}{-}{-} 3단계: 추상화 {-}{-}{-}}
\CommentTok{\# 학생 수가 5명이든 50명이든 같은 코드로 동작하도록 설계}

\CommentTok{\# {-}{-}{-} 4단계: 알고리즘 구현 {-}{-}{-}}
\NormalTok{best\_score }\OperatorTok{=}\NormalTok{ scores[}\DecValTok{0}\NormalTok{]      }\CommentTok{\# 첫 번째 점수를 임시 최고로 설정}
\NormalTok{best\_name  }\OperatorTok{=}\NormalTok{ students[}\DecValTok{0}\NormalTok{]}

\ControlFlowTok{for}\NormalTok{ i }\KeywordTok{in} \BuiltInTok{range}\NormalTok{(}\BuiltInTok{len}\NormalTok{(students)):}
    \ControlFlowTok{if}\NormalTok{ scores[i] }\OperatorTok{\textgreater{}}\NormalTok{ best\_score:}
\NormalTok{        best\_score }\OperatorTok{=}\NormalTok{ scores[i]}
\NormalTok{        best\_name  }\OperatorTok{=}\NormalTok{ students[i]}

\BuiltInTok{print}\NormalTok{(}\SpecialStringTok{f"최고 점수 학생: }\SpecialCharTok{\{}\NormalTok{best\_name}\SpecialCharTok{\}}\SpecialStringTok{ (}\SpecialCharTok{\{}\NormalTok{best\_score}\SpecialCharTok{\}}\SpecialStringTok{점)"}\NormalTok{)}
\end{Highlighting}
\end{Shaded}

\begin{tcolorbox}[enhanced jigsaw, arc=.35mm, bottomrule=.15mm, bottomtitle=1mm, breakable, colback=white, colbacktitle=quarto-callout-tip-color!10!white, colframe=quarto-callout-tip-color-frame, coltitle=black, left=2mm, leftrule=.75mm, opacityback=0, opacitybacktitle=0.6, rightrule=.15mm, title=\textcolor{quarto-callout-tip-color}{\faLightbulb}\hspace{0.5em}{파이썬 내장 함수로 더 간단하게}, titlerule=0mm, toprule=.15mm, toptitle=1mm]

위 반복문을 파이썬 내장 함수로 한 줄로 줄일 수 있습니다.

\begin{Shaded}
\begin{Highlighting}[]
\NormalTok{best\_idx   }\OperatorTok{=}\NormalTok{ scores.index(}\BuiltInTok{max}\NormalTok{(scores))}
\NormalTok{best\_name  }\OperatorTok{=}\NormalTok{ students[best\_idx]}
\NormalTok{best\_score }\OperatorTok{=}\NormalTok{ scores[best\_idx]}
\BuiltInTok{print}\NormalTok{(}\SpecialStringTok{f"최고 점수 학생: }\SpecialCharTok{\{}\NormalTok{best\_name}\SpecialCharTok{\}}\SpecialStringTok{ (}\SpecialCharTok{\{}\NormalTok{best\_score}\SpecialCharTok{\}}\SpecialStringTok{점)"}\NormalTok{)}
\end{Highlighting}
\end{Shaded}

두 방법의 결과는 같지만, 처음에는 직접 반복문을 써보며 알고리즘 흐름을
익히는 것이 중요합니다.

\end{tcolorbox}

\subsection{예제 2: 순서도를 코드로
옮기기}\label{uxc608uxc81c-2-uxc21cuxc11cuxb3c4uxb97c-uxcf54uxb4dcuxb85c-uxc62euxae30uxae30}

강의 본문의 ``평균 초과 학생 수'' 순서도를 파이썬으로 구현합니다.

\begin{Shaded}
\begin{Highlighting}[]
\NormalTok{scores }\OperatorTok{=}\NormalTok{ [}\DecValTok{78}\NormalTok{, }\DecValTok{91}\NormalTok{, }\DecValTok{65}\NormalTok{, }\DecValTok{82}\NormalTok{, }\DecValTok{75}\NormalTok{, }\DecValTok{88}\NormalTok{, }\DecValTok{70}\NormalTok{, }\DecValTok{95}\NormalTok{, }\DecValTok{60}\NormalTok{, }\DecValTok{83}\NormalTok{,}
          \DecValTok{77}\NormalTok{, }\DecValTok{89}\NormalTok{, }\DecValTok{72}\NormalTok{, }\DecValTok{68}\NormalTok{, }\DecValTok{91}\NormalTok{, }\DecValTok{84}\NormalTok{, }\DecValTok{76}\NormalTok{, }\DecValTok{80}\NormalTok{, }\DecValTok{55}\NormalTok{, }\DecValTok{93}\NormalTok{,}
          \DecValTok{74}\NormalTok{, }\DecValTok{87}\NormalTok{, }\DecValTok{69}\NormalTok{, }\DecValTok{85}\NormalTok{, }\DecValTok{78}\NormalTok{, }\DecValTok{92}\NormalTok{, }\DecValTok{63}\NormalTok{, }\DecValTok{79}\NormalTok{, }\DecValTok{88}\NormalTok{, }\DecValTok{71}\NormalTok{]}

\CommentTok{\# 평균 계산}
\NormalTok{avg }\OperatorTok{=} \BuiltInTok{sum}\NormalTok{(scores) }\OperatorTok{/} \BuiltInTok{len}\NormalTok{(scores)}
\BuiltInTok{print}\NormalTok{(}\SpecialStringTok{f"반 전체 평균: }\SpecialCharTok{\{}\NormalTok{avg}\SpecialCharTok{:.2f\}}\SpecialStringTok{점"}\NormalTok{)}

\CommentTok{\# 평균 초과 학생 수}
\NormalTok{count }\OperatorTok{=} \DecValTok{0}
\ControlFlowTok{for}\NormalTok{ s }\KeywordTok{in}\NormalTok{ scores:}
    \ControlFlowTok{if}\NormalTok{ s }\OperatorTok{\textgreater{}}\NormalTok{ avg:}
\NormalTok{        count }\OperatorTok{+=} \DecValTok{1}

\NormalTok{total }\OperatorTok{=} \BuiltInTok{len}\NormalTok{(scores)}
\BuiltInTok{print}\NormalTok{(}\SpecialStringTok{f"평균 초과 학생 수: }\SpecialCharTok{\{}\NormalTok{count}\SpecialCharTok{\}}\SpecialStringTok{명 / }\SpecialCharTok{\{}\NormalTok{total}\SpecialCharTok{\}}\SpecialStringTok{명"}\NormalTok{)}
\BuiltInTok{print}\NormalTok{(}\SpecialStringTok{f"비율: }\SpecialCharTok{\{}\NormalTok{count}\OperatorTok{/}\NormalTok{total}\OperatorTok{*}\DecValTok{100}\SpecialCharTok{:.1f\}}\SpecialStringTok{\%"}\NormalTok{)}
\end{Highlighting}
\end{Shaded}

\section{제출용 문제}\label{uxc81cuxcd9cuxc6a9-uxbb38uxc81c}

다음 요구사항을 만족하도록 \texttt{\#\ TODO}를 채워 \textbf{하나의
파이썬 파일(.py)}로 제출합니다.

\begin{center}\rule{0.5\linewidth}{0.5pt}\end{center}

\textbf{문제 1. 컴퓨팅 사고 4단계 적용 --- ``편의점 계산기''}

편의점에서 여러 상품을 구매할 때 총 금액을 계산하는 프로그램을 만드시오.
아래 데이터를 사용하고, CT 4단계를 주석으로 명시하시오.

\begin{Shaded}
\begin{Highlighting}[]
\NormalTok{items  }\OperatorTok{=}\NormalTok{ [}\StringTok{"아메리카노"}\NormalTok{, }\StringTok{"삼각김밥"}\NormalTok{, }\StringTok{"과자"}\NormalTok{, }\StringTok{"음료수"}\NormalTok{, }\StringTok{"껌"}\NormalTok{]}
\NormalTok{prices }\OperatorTok{=}\NormalTok{ [}\DecValTok{2000}\NormalTok{, }\DecValTok{1200}\NormalTok{, }\DecValTok{1500}\NormalTok{, }\DecValTok{1800}\NormalTok{, }\DecValTok{500}\NormalTok{]}
\NormalTok{qtys   }\OperatorTok{=}\NormalTok{ [}\DecValTok{2}\NormalTok{, }\DecValTok{3}\NormalTok{, }\DecValTok{1}\NormalTok{, }\DecValTok{2}\NormalTok{, }\DecValTok{4}\NormalTok{]   }\CommentTok{\# 각 상품의 구매 수량}

\CommentTok{\# {-}{-}{-} 1단계: 문제 분해 {-}{-}{-}}
\CommentTok{\# (작성하시오: 이 문제를 몇 개의 작은 문제로 나눌 수 있는가?)}

\CommentTok{\# {-}{-}{-} 2단계: 패턴 인식 {-}{-}{-}}
\CommentTok{\# (작성하시오: 어떤 패턴이 보이는가? 어떤 파이썬 구조를 쓸 것인가?)}

\CommentTok{\# {-}{-}{-} 3단계: 추상화 {-}{-}{-}}
\CommentTok{\# (작성하시오: 상품 수가 달라져도 동작하게 하려면 무엇을 일반화해야 하는가?)}

\CommentTok{\# {-}{-}{-} 4단계: 알고리즘 구현 {-}{-}{-}}
\NormalTok{total }\OperatorTok{=} \DecValTok{0}
\ControlFlowTok{for}\NormalTok{ i }\KeywordTok{in} \BuiltInTok{range}\NormalTok{(}\BuiltInTok{len}\NormalTok{(items)):}
    \CommentTok{\# }\AlertTok{TODO}\CommentTok{: 각 상품의 소계(prices[i] * qtys[i])를 total에 더하시오}
    \ControlFlowTok{pass}

\BuiltInTok{print}\NormalTok{(}\SpecialStringTok{f"총 결제 금액: }\SpecialCharTok{\{}\NormalTok{total}\SpecialCharTok{:,\}}\SpecialStringTok{원"}\NormalTok{)}
\end{Highlighting}
\end{Shaded}

\begin{center}\rule{0.5\linewidth}{0.5pt}\end{center}

\textbf{문제 2. 알고리즘 효율성 비교}

1부터 N까지의 \textbf{홀수}의 합을 구하는 두 가지 알고리즘을 완성하고,
결과가 같음을 확인하시오.

\begin{Shaded}
\begin{Highlighting}[]
\NormalTok{N }\OperatorTok{=} \DecValTok{100}

\CommentTok{\# 방법 1: 반복문으로 하나씩 확인}
\NormalTok{total\_loop }\OperatorTok{=} \DecValTok{0}
\ControlFlowTok{for}\NormalTok{ i }\KeywordTok{in} \BuiltInTok{range}\NormalTok{(}\DecValTok{1}\NormalTok{, N }\OperatorTok{+} \DecValTok{1}\NormalTok{):}
    \CommentTok{\# }\AlertTok{TODO}\CommentTok{: i가 홀수일 때만 total\_loop에 더하시오 (힌트: i \% 2 != 0)}
    \ControlFlowTok{pass}

\CommentTok{\# 방법 2: 수학 공식 사용}
\CommentTok{\# 힌트: 1부터 N까지 홀수의 개수는 N//2 이고,}
\CommentTok{\#       합은 (홀수 개수)\^{}2 입니다 (N=100이면 50\^{}2=2500)}
\NormalTok{total\_formula }\OperatorTok{=} \DecValTok{0}   \CommentTok{\# }\AlertTok{TODO}\CommentTok{: 수식으로 바로 계산하시오}

\BuiltInTok{print}\NormalTok{(}\SpecialStringTok{f"방법 1 (반복문): }\SpecialCharTok{\{}\NormalTok{total\_loop}\SpecialCharTok{\}}\SpecialStringTok{"}\NormalTok{)}
\BuiltInTok{print}\NormalTok{(}\SpecialStringTok{f"방법 2 (공식):   }\SpecialCharTok{\{}\NormalTok{total\_formula}\SpecialCharTok{\}}\SpecialStringTok{"}\NormalTok{)}
\BuiltInTok{print}\NormalTok{(}\SpecialStringTok{f"결과 동일 여부: }\SpecialCharTok{\{}\NormalTok{total\_loop }\OperatorTok{==}\NormalTok{ total\_formula}\SpecialCharTok{\}}\SpecialStringTok{"}\NormalTok{)}
\end{Highlighting}
\end{Shaded}

\begin{center}\rule{0.5\linewidth}{0.5pt}\end{center}

\textbf{문제 3. 나만의 알고리즘 설계 (서술형)}

일상생활에서 경험할 수 있는 문제 한 가지를 직접 골라, 아래 양식에 따라
작성하시오. (코드 작성은 선택 사항이며, 의사코드나 순서도 설명으로도
충분합니다.)

\begin{verbatim}
선택한 문제:

1단계 - 문제 분해 (작은 문제 3개 이상 나열):
  - 
  - 
  - 

2단계 - 패턴 인식 (반복 또는 조건 패턴이 있는가?):

3단계 - 추상화 (어떤 요소를 일반화할 수 있는가?):

4단계 - 알고리즘 (의사코드 또는 순서도 설명):
\end{verbatim}

\bookmarksetup{startatroot}

\chapter{1주차 연습문제}\label{sec-ch1-quiz}

다음 20문항은 4지선다입니다. 정답과 간단한 해설은 각 문항 아래 접힌
상자에 있습니다.

\begin{center}\rule{0.5\linewidth}{0.5pt}\end{center}

\textbf{1.} Jeannette Wing(2006)이 정의한 컴퓨팅 사고의 설명으로 가장
적절한 것은?

① 컴퓨터 언어를 유창하게 사용하는 능력\\
② 문제를 컴퓨터가 실행 가능한 방식으로 공식화하고 해결책을 표현하는 사고
과정\\
③ 프로그래밍 언어의 문법을 암기하는 것\\
④ 하드웨어 장치를 직접 조립하고 수리하는 기술

\begin{tcolorbox}[enhanced jigsaw, arc=.35mm, bottomrule=.15mm, bottomtitle=1mm, breakable, colback=white, colbacktitle=quarto-callout-note-color!10!white, colframe=quarto-callout-note-color-frame, coltitle=black, left=2mm, leftrule=.75mm, opacityback=0, opacitybacktitle=0.6, rightrule=.15mm, title=\textcolor{quarto-callout-note-color}{\faInfo}\hspace{0.5em}{노트}, titlerule=0mm, toprule=.15mm, toptitle=1mm]

정답 ②. Wing(2006)은 컴퓨팅 사고를 ``문제를 컴퓨터(또는 사람)가
효과적으로 실행할 수 있는 방식으로 공식화하고 그 해결책을 표현하는 사고
과정''으로 정의했습니다.

\end{tcolorbox}

\begin{center}\rule{0.5\linewidth}{0.5pt}\end{center}

\textbf{2.} 컴퓨팅 사고와 프로그래밍의 관계를 올바르게 설명한 것은?

① 컴퓨팅 사고와 프로그래밍은 완전히 같은 개념이다\\
② 프로그래밍 없이는 컴퓨팅 사고를 연습할 수 없다\\
③ 컴퓨팅 사고는 사고 과정이고, 프로그래밍은 그것을 컴퓨터에 전달하는
수단이다\\
④ 프로그래밍을 잘 하면 컴퓨팅 사고는 필요 없다

\begin{tcolorbox}[enhanced jigsaw, arc=.35mm, bottomrule=.15mm, bottomtitle=1mm, breakable, colback=white, colbacktitle=quarto-callout-note-color!10!white, colframe=quarto-callout-note-color-frame, coltitle=black, left=2mm, leftrule=.75mm, opacityback=0, opacitybacktitle=0.6, rightrule=.15mm, title=\textcolor{quarto-callout-note-color}{\faInfo}\hspace{0.5em}{노트}, titlerule=0mm, toprule=.15mm, toptitle=1mm]

정답 ③. 컴퓨팅 사고는 ``어떻게 풀 것인가''를 정리하는 사고 과정이며,
프로그래밍은 그 생각을 코드로 표현하는 수단입니다. 둘은 구분됩니다.

\end{tcolorbox}

\begin{center}\rule{0.5\linewidth}{0.5pt}\end{center}

\textbf{3.} 컴퓨터를 구성하는 두 가지 큰 요소는?

① 입력장치와 출력장치\\
② 하드웨어와 소프트웨어\\
③ CPU와 메모리\\
④ 운영체제와 응용프로그램

\begin{tcolorbox}[enhanced jigsaw, arc=.35mm, bottomrule=.15mm, bottomtitle=1mm, breakable, colback=white, colbacktitle=quarto-callout-note-color!10!white, colframe=quarto-callout-note-color-frame, coltitle=black, left=2mm, leftrule=.75mm, opacityback=0, opacitybacktitle=0.6, rightrule=.15mm, title=\textcolor{quarto-callout-note-color}{\faInfo}\hspace{0.5em}{노트}, titlerule=0mm, toprule=.15mm, toptitle=1mm]

정답 ②. 컴퓨터는 물리적 장치인 하드웨어와 하드웨어를 작동시키는 명령어
집합인 소프트웨어로 구성됩니다.

\end{tcolorbox}

\begin{center}\rule{0.5\linewidth}{0.5pt}\end{center}

\textbf{4.} 고급 프로그래밍 언어의 특징으로 옳은 것은?

① 컴퓨터가 바로 이해하는 0과 1로 이루어짐\\
② 기계어보다 실행 속도가 항상 빠름\\
③ 사람의 언어와 유사하며 배우기 쉽고 유지보수가 편리함\\
④ 하드웨어 장치를 직접 제어하는 데 가장 적합함

\begin{tcolorbox}[enhanced jigsaw, arc=.35mm, bottomrule=.15mm, bottomtitle=1mm, breakable, colback=white, colbacktitle=quarto-callout-note-color!10!white, colframe=quarto-callout-note-color-frame, coltitle=black, left=2mm, leftrule=.75mm, opacityback=0, opacitybacktitle=0.6, rightrule=.15mm, title=\textcolor{quarto-callout-note-color}{\faInfo}\hspace{0.5em}{노트}, titlerule=0mm, toprule=.15mm, toptitle=1mm]

정답 ③. 고급 언어(C, Java, Python 등)는 인간의 언어에 가까워 가독성이
높고 배우기 쉬우며 유지보수가 용이합니다.

\end{tcolorbox}

\begin{center}\rule{0.5\linewidth}{0.5pt}\end{center}

\textbf{5.} 소프트웨어 개발 6단계를 올바른 순서로 나열한 것은?

① 설계 → 문제 분석 → 프로그래밍 → 테스트 → 사용 → 유지보수\\
② 문제 분석 → 설계 → 프로그래밍 → 테스트 → 사용 → 유지보수\\
③ 프로그래밍 → 설계 → 문제 분석 → 테스트 → 유지보수 → 사용\\
④ 문제 분석 → 프로그래밍 → 설계 → 테스트 → 사용 → 유지보수

\begin{tcolorbox}[enhanced jigsaw, arc=.35mm, bottomrule=.15mm, bottomtitle=1mm, breakable, colback=white, colbacktitle=quarto-callout-note-color!10!white, colframe=quarto-callout-note-color-frame, coltitle=black, left=2mm, leftrule=.75mm, opacityback=0, opacitybacktitle=0.6, rightrule=.15mm, title=\textcolor{quarto-callout-note-color}{\faInfo}\hspace{0.5em}{노트}, titlerule=0mm, toprule=.15mm, toptitle=1mm]

정답 ②. 올바른 순서는 문제 분석(무엇을) → 설계(어떻게) →
프로그래밍(구현) → 테스트(검증) → 사용(배포) → 유지보수(개선)입니다.

\end{tcolorbox}

\begin{center}\rule{0.5\linewidth}{0.5pt}\end{center}

\textbf{6.} 소프트웨어 개발 단계 중 컴퓨팅 사고가 가장 직접적으로
활용되는 단계는?

① 1단계: 문제 분석\\
② 2단계: 설계\\
③ 4단계: 테스트\\
④ 6단계: 유지보수

\begin{tcolorbox}[enhanced jigsaw, arc=.35mm, bottomrule=.15mm, bottomtitle=1mm, breakable, colback=white, colbacktitle=quarto-callout-note-color!10!white, colframe=quarto-callout-note-color-frame, coltitle=black, left=2mm, leftrule=.75mm, opacityback=0, opacitybacktitle=0.6, rightrule=.15mm, title=\textcolor{quarto-callout-note-color}{\faInfo}\hspace{0.5em}{노트}, titlerule=0mm, toprule=.15mm, toptitle=1mm]

정답 ②. 설계 단계에서 알고리즘을 구체화하고 문제를 쪼개는 과정이 컴퓨팅
사고의 핵심 적용 단계입니다.

\end{tcolorbox}

\begin{center}\rule{0.5\linewidth}{0.5pt}\end{center}

\textbf{7.} 컴퓨팅 사고의 4단계 중 ``큰 문제를 작은 문제로 나누는
것''은?

① 패턴 인식\\
② 추상화\\
③ 문제 분해\\
④ 알고리즘

\begin{tcolorbox}[enhanced jigsaw, arc=.35mm, bottomrule=.15mm, bottomtitle=1mm, breakable, colback=white, colbacktitle=quarto-callout-note-color!10!white, colframe=quarto-callout-note-color-frame, coltitle=black, left=2mm, leftrule=.75mm, opacityback=0, opacitybacktitle=0.6, rightrule=.15mm, title=\textcolor{quarto-callout-note-color}{\faInfo}\hspace{0.5em}{노트}, titlerule=0mm, toprule=.15mm, toptitle=1mm]

정답 ③. 문제 분해(Decomposition)는 복잡한 문제를 관리 가능한 작은 부분
문제로 나누는 단계입니다.

\end{tcolorbox}

\begin{center}\rule{0.5\linewidth}{0.5pt}\end{center}

\textbf{8.} 마트에서 여러 물건을 살 때, 각각 ``(수량) × (단가)''를 반복
계산한다는 공통 규칙을 발견하는 사고 단계는?

① 문제 분해\\
② 패턴 인식\\
③ 추상화\\
④ 알고리즘

\begin{tcolorbox}[enhanced jigsaw, arc=.35mm, bottomrule=.15mm, bottomtitle=1mm, breakable, colback=white, colbacktitle=quarto-callout-note-color!10!white, colframe=quarto-callout-note-color-frame, coltitle=black, left=2mm, leftrule=.75mm, opacityback=0, opacitybacktitle=0.6, rightrule=.15mm, title=\textcolor{quarto-callout-note-color}{\faInfo}\hspace{0.5em}{노트}, titlerule=0mm, toprule=.15mm, toptitle=1mm]

정답 ②. 반복되는 공통 규칙이나 구조를 찾아내는 것이 패턴 인식(Pattern
Recognition)입니다.

\end{tcolorbox}

\begin{center}\rule{0.5\linewidth}{0.5pt}\end{center}

\textbf{9.} ``총 가격 = Σ (수량ᵢ × 단가ᵢ)''처럼 구체적인 상황에서
일반화된 공식을 만드는 단계는?

① 문제 분해\\
② 패턴 인식\\
③ 추상화\\
④ 알고리즘

\begin{tcolorbox}[enhanced jigsaw, arc=.35mm, bottomrule=.15mm, bottomtitle=1mm, breakable, colback=white, colbacktitle=quarto-callout-note-color!10!white, colframe=quarto-callout-note-color-frame, coltitle=black, left=2mm, leftrule=.75mm, opacityback=0, opacitybacktitle=0.6, rightrule=.15mm, title=\textcolor{quarto-callout-note-color}{\faInfo}\hspace{0.5em}{노트}, titlerule=0mm, toprule=.15mm, toptitle=1mm]

정답 ③. 추상화(Abstraction)는 불필요한 세부 사항을 제거하고 핵심만 남겨
일반적인 모델을 만드는 단계입니다.

\end{tcolorbox}

\begin{center}\rule{0.5\linewidth}{0.5pt}\end{center}

\textbf{10.} 알고리즘(Algorithm)의 정의로 가장 적절한 것은?

① 프로그래밍 언어의 문법 규칙\\
② 컴퓨터 하드웨어 구동 방식\\
③ 특정 문제를 해결하기 위한 명확하고 순서가 있는 명령어의 집합\\
④ 데이터를 저장하는 자료구조

\begin{tcolorbox}[enhanced jigsaw, arc=.35mm, bottomrule=.15mm, bottomtitle=1mm, breakable, colback=white, colbacktitle=quarto-callout-note-color!10!white, colframe=quarto-callout-note-color-frame, coltitle=black, left=2mm, leftrule=.75mm, opacityback=0, opacitybacktitle=0.6, rightrule=.15mm, title=\textcolor{quarto-callout-note-color}{\faInfo}\hspace{0.5em}{노트}, titlerule=0mm, toprule=.15mm, toptitle=1mm]

정답 ③. 알고리즘은 문제 해결을 위한 명확하고 유한하며 순서가 있는 절차
또는 명령어의 집합입니다.

\end{tcolorbox}

\begin{center}\rule{0.5\linewidth}{0.5pt}\end{center}

\textbf{11.} 좋은 알고리즘이 갖춰야 할 요건으로 적절하지 \textbf{않은}
것은?

① 각 단계가 명확해야 한다 (명확성)\\
② 반드시 끝나야 한다 (유한성)\\
③ 같은 결과를 내는 것 중 더 자원을 많이 쓰는 것이 좋다 (비효율성)\\
④ 같은 결과를 내는 것 중 더 빠른 것이 좋다 (효율성)

\begin{tcolorbox}[enhanced jigsaw, arc=.35mm, bottomrule=.15mm, bottomtitle=1mm, breakable, colback=white, colbacktitle=quarto-callout-note-color!10!white, colframe=quarto-callout-note-color-frame, coltitle=black, left=2mm, leftrule=.75mm, opacityback=0, opacitybacktitle=0.6, rightrule=.15mm, title=\textcolor{quarto-callout-note-color}{\faInfo}\hspace{0.5em}{노트}, titlerule=0mm, toprule=.15mm, toptitle=1mm]

정답 ③. 좋은 알고리즘은 명확성·유한성·효율성을 갖춰야 합니다. 자원을 더
많이 쓰는 것은 좋은 알고리즘의 요건이 아닙니다.

\end{tcolorbox}

\begin{center}\rule{0.5\linewidth}{0.5pt}\end{center}

\textbf{12.} 순서도(flowchart)에서 \textbf{조건 판단}을 나타내는 도형은?

① 타원\\
② 직사각형\\
③ 마름모\\
④ 평행사변형

\begin{tcolorbox}[enhanced jigsaw, arc=.35mm, bottomrule=.15mm, bottomtitle=1mm, breakable, colback=white, colbacktitle=quarto-callout-note-color!10!white, colframe=quarto-callout-note-color-frame, coltitle=black, left=2mm, leftrule=.75mm, opacityback=0, opacitybacktitle=0.6, rightrule=.15mm, title=\textcolor{quarto-callout-note-color}{\faInfo}\hspace{0.5em}{노트}, titlerule=0mm, toprule=.15mm, toptitle=1mm]

정답 ③. 마름모(다이아몬드)는 ``예/아니오''로 분기하는 조건 판단을
나타냅니다.

\end{tcolorbox}

\begin{center}\rule{0.5\linewidth}{0.5pt}\end{center}

\textbf{13.} 순서도에서 \textbf{시작 또는 끝}을 나타내는 도형은?

① 직사각형\\
② 마름모\\
③ 평행사변형\\
④ 타원(둥근 직사각형)

\begin{tcolorbox}[enhanced jigsaw, arc=.35mm, bottomrule=.15mm, bottomtitle=1mm, breakable, colback=white, colbacktitle=quarto-callout-note-color!10!white, colframe=quarto-callout-note-color-frame, coltitle=black, left=2mm, leftrule=.75mm, opacityback=0, opacitybacktitle=0.6, rightrule=.15mm, title=\textcolor{quarto-callout-note-color}{\faInfo}\hspace{0.5em}{노트}, titlerule=0mm, toprule=.15mm, toptitle=1mm]

정답 ④. 타원(또는 양쪽이 둥근 직사각형)은 순서도에서 시작 또는 종료를
나타냅니다.

\end{tcolorbox}

\begin{center}\rule{0.5\linewidth}{0.5pt}\end{center}

\textbf{14.} 다음 파이썬 코드에서 컴퓨팅 사고의 ``패턴 인식''을 직접
구현한 부분은?

\begin{Shaded}
\begin{Highlighting}[]
\NormalTok{scores }\OperatorTok{=}\NormalTok{ [}\DecValTok{85}\NormalTok{, }\DecValTok{92}\NormalTok{, }\DecValTok{78}\NormalTok{, }\DecValTok{96}\NormalTok{, }\DecValTok{88}\NormalTok{]}
\NormalTok{count }\OperatorTok{=} \DecValTok{0}
\ControlFlowTok{for}\NormalTok{ s }\KeywordTok{in}\NormalTok{ scores:}
    \ControlFlowTok{if}\NormalTok{ s }\OperatorTok{\textgreater{}} \DecValTok{90}\NormalTok{:}
\NormalTok{        count }\OperatorTok{+=} \DecValTok{1}
\end{Highlighting}
\end{Shaded}

① \texttt{scores\ =\ {[}85,\ 92,\ 78,\ 96,\ 88{]}}\\
② \texttt{count\ =\ 0}\\
③ \texttt{for\ s\ in\ scores:} 와 \texttt{if\ s\ \textgreater{}\ 90:}\\
④ \texttt{count\ +=\ 1}

\begin{tcolorbox}[enhanced jigsaw, arc=.35mm, bottomrule=.15mm, bottomtitle=1mm, breakable, colback=white, colbacktitle=quarto-callout-note-color!10!white, colframe=quarto-callout-note-color-frame, coltitle=black, left=2mm, leftrule=.75mm, opacityback=0, opacitybacktitle=0.6, rightrule=.15mm, title=\textcolor{quarto-callout-note-color}{\faInfo}\hspace{0.5em}{노트}, titlerule=0mm, toprule=.15mm, toptitle=1mm]

정답 ③. 반복 패턴(\texttt{for})과 조건 패턴(\texttt{if})이 컴퓨팅 사고의
패턴 인식을 코드로 표현한 것입니다.

\end{tcolorbox}

\begin{center}\rule{0.5\linewidth}{0.5pt}\end{center}

\textbf{15.} 1부터 100까지의 합을 구하는 두 방법 중 더 효율적인 것은?

① \texttt{total\ =\ 0;\ for\ i\ in\ range(1,\ 101):\ total\ +=\ i}\\
② \texttt{total\ =\ 100\ *\ 101\ //\ 2}\\
③ 두 방법의 효율은 완전히 같다\\
④ 효율성은 결과가 다를 때만 비교할 수 있다

\begin{tcolorbox}[enhanced jigsaw, arc=.35mm, bottomrule=.15mm, bottomtitle=1mm, breakable, colback=white, colbacktitle=quarto-callout-note-color!10!white, colframe=quarto-callout-note-color-frame, coltitle=black, left=2mm, leftrule=.75mm, opacityback=0, opacitybacktitle=0.6, rightrule=.15mm, title=\textcolor{quarto-callout-note-color}{\faInfo}\hspace{0.5em}{노트}, titlerule=0mm, toprule=.15mm, toptitle=1mm]

정답 ②. 방법 2는 수학 공식을 사용해 단 한 번의 계산으로 끝나므로, N이
아무리 커져도 같은 속도입니다. 방법 1은 N번 반복해야 합니다.

\end{tcolorbox}

\begin{center}\rule{0.5\linewidth}{0.5pt}\end{center}

\textbf{16.} ``함수(function)''는 컴퓨팅 사고의 어느 단계와 가장
직접적으로 관련되는가?

① 문제 분해\\
② 패턴 인식\\
③ 추상화\\
④ 알고리즘

\begin{tcolorbox}[enhanced jigsaw, arc=.35mm, bottomrule=.15mm, bottomtitle=1mm, breakable, colback=white, colbacktitle=quarto-callout-note-color!10!white, colframe=quarto-callout-note-color-frame, coltitle=black, left=2mm, leftrule=.75mm, opacityback=0, opacitybacktitle=0.6, rightrule=.15mm, title=\textcolor{quarto-callout-note-color}{\faInfo}\hspace{0.5em}{노트}, titlerule=0mm, toprule=.15mm, toptitle=1mm]

정답 ③. 함수는 복잡한 처리를 하나의 이름으로 묶어 세부 사항을 숨기는
추상화의 대표적 도구입니다.

\end{tcolorbox}

\begin{center}\rule{0.5\linewidth}{0.5pt}\end{center}

\textbf{17.} 다음 중 알고리즘을 표현하는 방법이 \textbf{아닌} 것은?

① 자연어(우리말 설명)\\
② 순서도(Flowchart)\\
③ 의사코드(Pseudocode)\\
④ 컴파일러(Compiler)

\begin{tcolorbox}[enhanced jigsaw, arc=.35mm, bottomrule=.15mm, bottomtitle=1mm, breakable, colback=white, colbacktitle=quarto-callout-note-color!10!white, colframe=quarto-callout-note-color-frame, coltitle=black, left=2mm, leftrule=.75mm, opacityback=0, opacitybacktitle=0.6, rightrule=.15mm, title=\textcolor{quarto-callout-note-color}{\faInfo}\hspace{0.5em}{노트}, titlerule=0mm, toprule=.15mm, toptitle=1mm]

정답 ④. 컴파일러는 소스코드를 실행 가능한 형태로 변환하는 도구이며,
알고리즘 표현 방법이 아닙니다. 알고리즘은 자연어·순서도·의사코드로
표현합니다.

\end{tcolorbox}

\begin{center}\rule{0.5\linewidth}{0.5pt}\end{center}

\textbf{18.} 소프트웨어 개발에서 ``테스트(4단계)''는 누가 수행하는 것이
바람직한가?

① 반드시 소프트웨어를 만든 개발자가 직접 수행해야 한다\\
② 개발자가 아닌 다른 사람이, 최대한 다양한 환경에서 수행해야 한다\\
③ 사용자 배포 후 실제 사용 과정에서 자동으로 수행된다\\
④ 테스트 단계는 소규모 프로젝트에서는 생략 가능하다

\begin{tcolorbox}[enhanced jigsaw, arc=.35mm, bottomrule=.15mm, bottomtitle=1mm, breakable, colback=white, colbacktitle=quarto-callout-note-color!10!white, colframe=quarto-callout-note-color-frame, coltitle=black, left=2mm, leftrule=.75mm, opacityback=0, opacitybacktitle=0.6, rightrule=.15mm, title=\textcolor{quarto-callout-note-color}{\faInfo}\hspace{0.5em}{노트}, titlerule=0mm, toprule=.15mm, toptitle=1mm]

정답 ②. 개발자 본인은 자신의 코드에 익숙해 실수를 놓치기 쉽습니다. 다른
사람이 다양한 환경에서 테스트해야 더 많은 오류를 발견할 수 있습니다.

\end{tcolorbox}

\begin{center}\rule{0.5\linewidth}{0.5pt}\end{center}

\textbf{19.} 파이썬이 ``인터프리터 방식'' 언어라는 의미는?

① 소스코드 전체를 한 번에 기계어로 변환한 뒤 실행한다\\
② 소스코드를 한 줄씩 즉시 해석하고 실행한다\\
③ 기계어로 작성된 코드를 사람이 읽을 수 있게 변환한다\\
④ 실행 전에 반드시 컴파일 과정이 필요하다

\begin{tcolorbox}[enhanced jigsaw, arc=.35mm, bottomrule=.15mm, bottomtitle=1mm, breakable, colback=white, colbacktitle=quarto-callout-note-color!10!white, colframe=quarto-callout-note-color-frame, coltitle=black, left=2mm, leftrule=.75mm, opacityback=0, opacitybacktitle=0.6, rightrule=.15mm, title=\textcolor{quarto-callout-note-color}{\faInfo}\hspace{0.5em}{노트}, titlerule=0mm, toprule=.15mm, toptitle=1mm]

정답 ②. 인터프리터(interpreter) 방식은 소스코드를 한 줄씩 즉시
해석·실행합니다. 파이썬이 이 방식을 사용하기 때문에 바로 실행 결과를
확인하기 편리합니다.

\end{tcolorbox}

\begin{center}\rule{0.5\linewidth}{0.5pt}\end{center}

\textbf{20.} 이 강의의 15주 구성에서 ``후반부(9\textasciitilde15주)''의
주요 주제는?

① 스크래치를 이용한 블록 코딩\\
② 웹 개발과 HTML/CSS 기초\\
③ 데이터 분석과 인공지능 기반 문제해결\\
④ 자바를 이용한 객체지향 프로그래밍

\begin{tcolorbox}[enhanced jigsaw, arc=.35mm, bottomrule=.15mm, bottomtitle=1mm, breakable, colback=white, colbacktitle=quarto-callout-note-color!10!white, colframe=quarto-callout-note-color-frame, coltitle=black, left=2mm, leftrule=.75mm, opacityback=0, opacitybacktitle=0.6, rightrule=.15mm, title=\textcolor{quarto-callout-note-color}{\faInfo}\hspace{0.5em}{노트}, titlerule=0mm, toprule=.15mm, toptitle=1mm]

정답 ③. 이 강의의 후반부(9\textasciitilde15주)는 함수·데이터
분석·시각화·패턴 탐색·AI 원리·생성형 AI·프로젝트로 구성됩니다.

\end{tcolorbox}




\end{document}
